% Demonstration of the unofficial Fusion CDT Beamer theme.
% All content is illustrative. Build from this folder with: latexmk
% SPDX-License-Identifier: CC0-1.0
\documentclass[aspectratio=169, 11pt]{beamer}

\usetheme{FusionCDT}

\usepackage{amsmath}
\usepackage{booktabs}
\usepackage{pgfplots}
\pgfplotsset{compat=1.18}
\usetikzlibrary{positioning}

\title{Presentation title}
\subtitle{Optional subtitle}
\author[A. Name]{Author Name}
\institute{Affiliation}
\date{\today}
\hypersetup{pdfsubject={Demonstration of an unofficial, community-maintained
  Beamer template inspired by the Fusion CDT visual identity}}

\begin{document}

\fcdttitleframe

\begin{frame}{Outline}
  \tableofcontents
\end{frame}

\section{Text and lists}

\begin{frame}{Lists}{Frames can have an optional subtitle}
  \begin{itemize}
    \item Bullet points use small blue squares.
      \begin{itemize}
        \item Second-level points use a magenta dash.
        \item Inline mathematics: \(\nabla \cdot \mathbf{B} = 0\).
      \end{itemize}
    \item Use \alert{alerted text} sparingly for key words.
  \end{itemize}
  \begin{enumerate}
    \item Numbered lists use blue squares.
    \item They work with overlays and all usual Beamer features.
  \end{enumerate}
\end{frame}

\begin{frame}{Blocks}
  \begin{columns}[T]
    \begin{column}{0.48\textwidth}
      \begin{block}{Block}
        For definitions and key results.
      \end{block}
      \begin{exampleblock}{Example block}
        For worked examples.
      \end{exampleblock}
    \end{column}
    \begin{column}{0.48\textwidth}
      \begin{alertblock}{Alert block}
        For warnings and take-home messages.
      \end{alertblock}
    \end{column}
  \end{columns}
\end{frame}

\section{Mathematics and figures}

\begin{frame}{Equations}{Standard definitions, shown as an example}
  The electron plasma frequency and the Debye length are
  \begin{align}
    \omega_{pe} &= \sqrt{\frac{n_e e^2}{\varepsilon_0 m_e}},
    &
    \lambda_D &= \sqrt{\frac{\varepsilon_0 k_B T_e}{n_e e^2}}.
  \end{align}
  Displayed equations are numbered, and inline mathematics such as
  \(\lambda_D \, \omega_{pe} = \sqrt{k_B T_e / m_e}\) sits in the text.
\end{frame}

\begin{frame}{Plots}{Illustrative functions, not data}
  \centering
  \begin{tikzpicture}
    \begin{axis}[
        width=0.92\textwidth, height=0.62\textheight,
        xlabel={\(x\)}, ylabel={\(f(x) = \exp\bigl(-x^2 / 2\sigma^2\bigr)\)},
        axis lines=left, grid=major, grid style={fcdtPaleBlue},
        legend pos=north east, legend cell align=left,
        legend style={draw=none, fill=none, font=\footnotesize},
        mark repeat=6, mark size=1.6pt, ymin=0, ymax=1.05]
      \addplot+[domain=-3:3, samples=61] {exp(-x^2/(2*0.5^2))};
      \addplot+[domain=-3:3, samples=61] {exp(-x^2/(2*0.8^2))};
      \addplot+[domain=-3:3, samples=61] {exp(-x^2/(2*1.1^2))};
      \addplot+[domain=-3:3, samples=61] {exp(-x^2/(2*1.4^2))};
      \legend{\(\sigma = 0.5\), \(\sigma = 0.8\), \(\sigma = 1.1\), \(\sigma = 1.4\)}
    \end{axis}
  \end{tikzpicture}
\end{frame}

\begin{frame}{Figures}
  \begin{columns}[c]
    \begin{column}{0.55\textwidth}
      \begin{figure}
        \centering
        \begin{tikzpicture}[
            box/.style={draw=fcdtBlue, fill=fcdtBlueTint, rounded corners=2pt,
              minimum width=2cm, minimum height=0.9cm, font=\small},
            arrow/.style={-latex, thick, fcdtDarkBlue}]
          \node[box] (a) {Input};
          \node[box, right=0.6cm of a] (b) {Model};
          \node[box, right=0.6cm of b] (c) {Output};
          \draw[arrow] (a) -- (b);
          \draw[arrow] (b) -- (c);
        \end{tikzpicture}
        \caption{An illustrative schematic drawn with TikZ.}
      \end{figure}
    \end{column}
    \begin{column}{0.4\textwidth}
      \begin{itemize}
        \item Figures and captions follow the theme colours.
        \item Use \texttt{\textbackslash{}includegraphics} for image files.
      \end{itemize}
    \end{column}
  \end{columns}
\end{frame}

\begin{frame}{Tables}{The main theme colours}
  \centering
  \begin{tabular}{@{}llll@{}}
    \toprule
    & Colour name & Hex value & Used for \\
    \midrule
    \textcolor{fcdtBlue}{\rule{1.2em}{1.2ex}}     & \texttt{fcdtBlue}     & \texttt{\#277CB3} & Blue slides, footer, block titles \\
    \textcolor{fcdtDarkBlue}{\rule{1.2em}{1.2ex}} & \texttt{fcdtDarkBlue} & \texttt{\#1D6394} & Frame titles, structure \\
    \textcolor{fcdtMagenta}{\rule{1.2em}{1.2ex}}  & \texttt{fcdtMagenta}  & \texttt{\#DC2099} & Accents, second plot series \\
    \textcolor{fcdtLime}{\rule{1.2em}{1.2ex}}     & \texttt{fcdtLime}     & \texttt{\#C5D300} & Accent stripe only \\
    \bottomrule
  \end{tabular}
\end{frame}

\section{References}

\begin{frame}{Citations}
  Citations work as usual. This theme builds on Beamer~\cite{beamer}, the
  plots use pgfplots~\cite{pgfplots}, and the typeface is
  Figtree~\cite{figtree}.
\end{frame}

\begin{frame}{References}
  \footnotesize
  \bibliographystyle{plain}
  \bibliography{references}
\end{frame}

\begin{frame}{About this template}
  \footnotesize
  Unofficial, community-maintained LaTeX Beamer template inspired by the
  Fusion CDT visual identity. This is not an official Fusion CDT template and
  is not endorsed by Fusion CDT or its partner institutions.

  \medskip
  The Fusion CDT logo belongs to its owners and is downloaded separately; it
  is not covered by the template's licence. The Figtree typeface is used under
  the SIL Open Font License~1.1.
\end{frame}

\fcdtclosingframe{Thank you}

\end{document}
